The unexecuted limit orders resting at a venue are collectively called the
\emph{limit order book}.
This section makes that object precise.

\begin{defi}
\label{def.limitOrderBook}
The \emph{limit order book} at time $t$ is the pair of functions
\begin{equation}
\label{eq.limitOrderBook}
p \mapsto \bidVolumePriceP, \qquad p \mapsto \askVolumePriceP, \qquad p \in \tickSizeOfLOB \Z,
\end{equation}
where $\bidVolumePriceP$ is the total size of the buy limit orders resting at price $p$ at
time $t$,
and $\askVolumePriceP$ is the total size of the sell limit orders resting at price $p$ at
time $t$.
\end{defi}

Both functions take non-negative integer values and vanish outside a finite set of prices.
It is convenient to picture the book as two queues facing one another across a gap:
the bid side, occupying the low prices,
and the ask side, occupying the high prices.

\begin{defi}
\label{def.bestQuotes}
The \emph{best bid price} and the \emph{best ask price} at time $t$ are
\begin{equation}
\label{eq.bestQuotes}
\bestBidPrice_t := \max \lbrace p : \bidVolumePriceP > 0 \rbrace,
\qquad
\bestAskPrice_t := \min \lbrace p : \askVolumePriceP > 0 \rbrace,
\end{equation}
and the volumes resting there,
\begin{equation*}
\bestBidVolume_t := \volume^{b}_{t}(\bestBidPrice_t),
\qquad
\bestAskVolume_t := \volume^{a}_{t}(\bestAskPrice_t),
\end{equation*}
are called the \emph{depth} at the best quotes.
Together, the four numbers $(\bestBidPrice_t, \bestBidVolume_t, \bestAskPrice_t, \bestAskVolume_t)$
constitute the \emph{top of the book}.
\end{defi}

The two sides never overlap.
If a buy limit order arrives priced at or above $\bestAskPrice_t$,
it does not rest in the book:
it is matched immediately against the resting sell orders,
exactly as a market order would be.
Consequently
\begin{equation}
\label{eq.positiveSpread}
\LOBspread_t := \bestAskPrice_t - \bestBidPrice_t \geq \tickSizeOfLOB
\end{equation}
at all times.
The quantity $\LOBspread_t$ is the \emph{bid-ask spread}.
It is the price of immediacy:
a trader who buys at the ask and immediately sells at the bid
pays $\LOBspread_t$ per unit for the privilege of not having waited.

Two summaries of the top of the book are used constantly in practice.
The first is the \emph{mid-price},
\begin{equation}
\label{eq.midPrice}
\midPrice_t := \frac{\bestBidPrice_t + \bestAskPrice_t}{2},
\end{equation}
which is the usual answer to the question ``what is the price?'',
even though it is a price at which nobody can trade.
The second is the \emph{queue imbalance},
\begin{equation}
\label{eq.volumeImbalance}
\volumeImbalance_t := \frac{\bestBidVolume_t - \bestAskVolume_t}{\bestBidVolume_t + \bestAskVolume_t} \in [-1,1],
\end{equation}
which measures how lopsided the top of the book is.
The imbalance is one of the most robust short-horizon predictors of the next mid-price move:
when the bid queue is much longer than the ask queue,
the ask queue tends to be exhausted first, and the price tends to rise.

\begin{remark}
\label{remark.midPriceIsNotTradeable}
Equation \eqref{eq.midPrice} defines a number, not a trading opportunity.
The gap between the mid-price and the price actually obtained
is where most of the cost of trading hides,
and it is the reason why a backtest that fills orders at the mid-price
flatters a strategy.
\end{remark}

\subsection{Execution priority}
\label{sec.executionPriority}

When a market order arrives, the venue must decide which resting limit orders it consumes.
Almost all venues apply \emph{price-time priority}:
better prices are served first,
and among orders at the same price,
the one that arrived earlier is served first.

\begin{defi}
\label{def.executionPriority}
Let $o$ and $o\derivative$ be two resting buy limit orders,
with prices $p, p\derivative$ and submission times $s, s\derivative$.
We say that $o$ has \emph{higher priority} than $o\derivative$ if
\begin{equation}
\label{eq.orderingOfLimitOrders}
p > p\derivative,
\qquad \text{ or } \qquad
p = p\derivative \, \text{ and } \, s < s\derivative.
\end{equation}
The same definition applies to sell limit orders with the first inequality reversed.
\end{defi}

Price-time priority is what makes a limit order a queueing problem.
An order submitted at the back of a long queue at the best bid
may wait a long time,
and it may never be executed at all if the price moves away.
The trader's position in the queue is therefore an asset,
and the decision of when to cancel and re-post is a real one:
cancelling forfeits the priority already earned.
